\begin{abstract}
Self-evolving agents convert previous trajectories into reusable textual memory. Existing analyses usually treat the retrieved text as the relevant state variable. We study a second variable carried into that state: the provenance of the trajectory from which the memory was written. A recent financial-agent audit reports a sharp ReasoningBank split. Success-derived retrievals achieve 0.931 accuracy over 101 cases, while failure-derived retrievals achieve 0.647 over 34 cases. All ten reported persistent W$\rightarrow$W failures occur in the failure-derived stratum. We formalize provenance as an upstream variable in memory construction and execute a controlled matched-memory intervention on released WebArena evidence. Six candidate query pairs are frozen before downstream outcomes. Four pass a pre-outcome information-equivalence gate. The gate requires semantic similarity. An independent verifier checks guidance equivalence. Literal benchmark-reference leakage is excluded. Across these four pairs, 48 downstream requests complete with zero support failures. Success-derived memories achieve mean terminal success 0.375, compared with 0.208 for matched failure-derived memories, a +0.1667 difference. This exceeds the preregistered 0.15 effect floor, while the frozen 100,000-permutation test gives $p=0.0785$ and therefore does not pass the $0.05$ support gate. No counterevidence gate fires. The controlled direction agrees with the persistence pattern in the financial audit while leaving the causal claim active pending stronger confirmation. These results motivate provenance-aware memory governance: future reuse decisions should retain evidence about how a memory was learned, not only the text that was stored.
\end{abstract}
